%%% ───────────────────────────────────────────────────────────────────────────
%%% §2 / §3  METHOD  (paste-ready LaTeX skeleton)
%%% ───────────────────────────────────────────────────────────────────────────

\section{Method}
\label{sec:method}

\subsection{Problem Setup}
\label{sec:setup}

We are given two token-aligned model variants that share a vocabulary and
tokenizer: an \emph{organism} model~$M_{\mathrm{org}}$ whose behavior we wish
to monitor, and a \emph{base} model~$M_{\mathrm{base}}$ that serves as a
distributional reference. In our running experiment, $M_{\mathrm{org}}$ is a
reward-model LoRA trained to exploit evaluation
biases~\citep{sheshadri2025replication, marks2025auditing}, and
$M_{\mathrm{base}}$ is the unmodified instruction-tuned checkpoint.

\paragraph{Notation.}
Let $d$ be the residual-stream dimension and $T$ the sequence length. We
write $h^{\mathrm{org}}_{i,t}\in\mathbb{R}^{d}$ for the hidden state of
$M_{\mathrm{org}}$ at token position~$t$ of the $i$-th example, and
$h^{\mathrm{base}}_{i,t}$ for the token-aligned hidden state of $M_{\mathrm{base}}$
run on the \emph{same} text (prefilled).

\paragraph{Annotated dataset.}
Let
\begin{equation}
  \mathcal{D} \;=\;
  \bigl\{\,(\text{prompt}_i,\;\text{response}_i,\;\tau_i)\bigr\}_{i=1}^{N},
\end{equation}
where $\tau_i\in\mathbb{N}$ is the \emph{onset token index}: the
response-relative position at which the model first commits to the target
behavior in example~$i$. We follow~\citet{snel2025first} in using the first
token of an annotated span as the onset, motivated by their empirical finding
that first-token AUROC dramatically exceeds in-span AUROC.

\paragraph{Representational basis.}
Let $\phi_B:\mathbb{R}^d\!\to\mathbb{R}^k$ be a \emph{basis map} that projects a
hidden state into a $k$-dimensional feature space. The framework is
indifferent to~$B$; concrete instantiations are discussed in
Section~\ref{sec:basis}.

\subsection{Activation Templates via Cohort Averaging}
\label{sec:templates}

\paragraph{Per-token features and centered delta.}
For each example~$i$ and token~$t$ we define the \emph{centered-delta
feature}
\begin{equation}
  \label{eq:centered-delta}
  x_{i,t}
  \;=\;
  \phi_B\!\bigl(h^{\mathrm{org}}_{i,t}\bigr)
  \;-\;
  \phi_B\!\bigl(h^{\mathrm{base}}_{i,t}\bigr)
  \;-\;
  \underbrace{%
    \frac{1}{T_i}\!\sum_{t'=1}^{T_i}\!
    \Bigl[\phi_B\!\bigl(h^{\mathrm{org}}_{i,t'}\bigr)
          -\phi_B\!\bigl(h^{\mathrm{base}}_{i,t'}\bigr)\Bigr]
  }_{\text{response-level mean, subtracted}}.
\end{equation}
Subtracting the model-pair difference isolates the organism-specific signal;
mean-centering removes content correlations shared with text. Ablation in
Section~\ref{sec:ablations} shows the centering recovers $\sim$14 percentage
points of localization performance over un-centered baselines.

\paragraph{Cohort-averaged template.}
Fix a half-window $W\in\mathbb{N}$. The \emph{activation template} (or
\emph{onset kernel}) is the cohort average of centered-delta windows aligned
to the annotated onset:
\begin{equation}
  \label{eq:template}
  K[w]
  \;=\;
  \frac{1}{N}\sum_{i=1}^{N} x_{i,\,\tau_i + w},
  \qquad w\in\{-W,\ldots,W-1\},
\end{equation}
yielding a matrix $K\in\mathbb{R}^{k\times 2W}$. Each column~$K[\cdot]$ is the
mean feature value at offset~$w$ from onset; $K$ is therefore a
\emph{prototype trajectory} of feature evolution around behavioral
commitment. We additionally smooth $K$ along the time axis with a
short rolling mean (window 9, boundary-respecting) to suppress
annotation-jitter noise, following~\citet{schroeter2021misaligned}.

\subsection{Detection via Cross-Correlation}
\label{sec:detection}

Given a new response with feature trajectory $\{x_t\}_{t=1}^{T}$, we compute
the \emph{detection score} at each candidate onset position~$t$ by sliding
the template:
\begin{equation}
  \label{eq:detection}
  s(t)
  \;=\;
  \sum_{w=-W}^{W-1}
  \bigl\langle K[w],\; x_{t+w} \bigr\rangle,
  \qquad
  t \in \{W+1,\ldots,T-W\},
\end{equation}
with predicted onset $\hat{\tau} = \arg\max_t s(t)$. We unit-normalize
each row of~$K$ before sliding, reducing Equation~\eqref{eq:detection} to
cosine correlation per feature.

\paragraph{Connection to matched filtering.}
Equation~\eqref{eq:detection} is a discrete cross-correlation between
template~$K$ and trajectory $\{x_t\}$ -- a temporal \emph{matched filter}
in feature space. Under the additive Gaussian-noise model
$x_t = K[t-\tau] + \varepsilon_t$ with covariance~$\Sigma$, the optimal
Neyman--Pearson decision statistic is the whitened inner product
\begin{equation}
  s^{*}(t) = \sum_{w} K[w]^\top \Sigma^{-1} x_{t+w},
\end{equation}
which reduces to Equation~\eqref{eq:detection} when~$\Sigma \propto I$
\citep{parra2005recipes, franke2015bayes}. We import this framework wholesale
from single-trial event-related potential (ERP) detection and
Bayes-optimal spike sorting; both share the same problem structure
(stereotyped waveform embedded in noise with unknown onset time) where
template cross-correlation is the maximum-likelihood detector.

\paragraph{Tolerance evaluation.}
Following the action-spotting literature~\citep{giancola2018soccernet}, we
score predictions using \emph{positional hit-rate at tolerance~$W$}: a
prediction $\hat{\tau}$ is correct if $|\hat{\tau} - \tau_{\mathrm{gt}}| \le W$
for any annotated onset $\tau_{\mathrm{gt}}$ in the response. This handles
inherent imprecision in behavioral onset annotation, well-characterized in
adjacent fields~\citep{schroeter2021misaligned} but novel for LLM activation
analysis.

\paragraph{Connection to DTW.}
We considered Soft-DTW~\citep{cuturi2017softdtw} for handling onset jitter
but rejected it: annotation variability in our dataset is small ($\sim$2--5
tokens median error), well below the regime where DTW alignment adds value,
and fixed-grid cross-correlation in Equation~\eqref{eq:detection} avoids
extra hyperparameters.

\subsection{Basis Agnosticism}
\label{sec:basis}

A design goal is that $\phi_B$ is a plug-in parameter, not an architectural
commitment. Equations~\eqref{eq:centered-delta}--\eqref{eq:detection} are
well-defined for any linear or nonlinear~$\phi_B$. We study five
instantiations:
\begin{enumerate}
  \item \textbf{Trait projections.}
    $\phi_B(h) = (v_1^\top h,\,\ldots,\,v_k^\top h)$ with $v_j$ a behavioral
    probe direction extracted via contrastive
    mean-difference~\citep{panickssery2024caa}. Our default uses 173
    emotion-set traits.
  \item \textbf{Raw residual stream (whitened).} $\phi_B(h) = W h$ where $W$ is
    a PCA whitening matrix; $k=d$.
  \item \textbf{Random projections.} $\phi_B(h) = \frac{1}{\sqrt{m}} R h$ with
    $R$ Gaussian. By Johnson--Lindenstrauss, inner products are approximately
    preserved. Calibration baseline.
  \item \textbf{PCA-of-delta.} Top-K principal components of the
    $\{(h^{\mathrm{org}}_{i,\tau_i} - h^{\mathrm{base}}_{i,\tau_i})\}_i$
    matrix at annotated onset tokens. Data-driven basis derived from
    annotated events.
  \item \textbf{LoRA-direction projection.} For each layer~$\ell$, projection
    onto the LoRA's $B_\ell$-matrix columns
    (rank~$r_\ell \in \{16, \ldots, 64\}$). The most principled basis: by
    construction, the directions through which the LoRA can intervene.
\end{enumerate}
The choice trades interpretability against task-specificity: trait
projections yield human-readable template rows
(``\textsf{shame} rises 4 tokens before onset''); raw residuals require no
annotation but produce opaque columns. We report results across all five.

\subsection{Per-Archetype Templates}
\label{sec:archetypes}

Not all behavioral commitments share temporal shape. We cluster training
examples by their cohort-averaged shape (Ward linkage on Frobenius cosine
similarity matrix between per-bias templates), recovering $k=3$ archetypes
in our reward-hacking corpus. For each cluster~$c$ with subset
$\mathcal{D}_c$, we fit a dedicated template $K^{(c)}$ via
Equation~\eqref{eq:template} restricted to $\mathcal{D}_c$. At detection
time, we max-pool over cluster scores:
\begin{equation}
  s^{\mathrm{ens}}(t) = \max_c\; s^{(c)}(t).
\end{equation}
Section~\ref{sec:results} shows that per-archetype templates outperform a
single universal template by~$\sim$5 percentage points held-out, while
single-bias templates overfit the training cohort.
